\section{Bounded detuning, interior splits, and residue classes}
\label{inf:sec-widened}

Fix $0<\varepsilon_{\mathrm{sp}}<1/4$ and write
\[
 I_{\mathrm{sp}}=[\varepsilon_{\mathrm{sp}}^2,
                         (1/2-\varepsilon_{\mathrm{sp}})^2],\qquad
 \varepsilon_{\mathrm{sp}}\le r\le1/2-\varepsilon_{\mathrm{sp}}.
\]
Constants in this section may depend on this interval and on a fixed
upper bound $T$ for $\tau$. The coefficient spaces and their angular
radius two remain those of Definition~\ref{inf:spaces}.

\begin{lemma}[Low-state transfer on a fixed wider window]
\label{inf:wide-low-transfer}
Fix $C_{\mathrm{ap}}>0$. The second conclusion of
Proposition~\ref{inf:low-exclusion} holds for
$|p-p_0|\le C_{\mathrm{ap}}/p_0$, after increasing the threshold
depending on $C_{\mathrm{ap}}$. More generally, its proof is uniform
for $C_{\mathrm{ap}}=O(\log p_0)$.
\end{lemma}
\begin{proof}
Suppose a nonseed low degree $\ell$ has a zero $z(p)$ with
$|z(p)^2-R(p)^2|\le32$. Continue it and the Neumann zero along the
interval to $p_0$. Equation~\eqref{inf:order-lipschitz} first gives
$|z(t)-R(t)|\le C(1+C_{\mathrm{ap}})/p_0$ there, so both radii are
comparable to $2p_0$ and both order-to-radius ratios belong to
$[1/3,2/3]$. Put
\[
 f_{\mathrm{ord}}(x)=\arccos x/\sqrt{1-x^2}.
\]
Its derivative is bounded on that compact interval.
Equation~\eqref{inf:order-derivative} therefore gives
\[
 |z'(t)-R'(t)|\le C(\ell+1)/p_0
                         +C(1+C_{\mathrm{ap}})/p_0^2.
\]
Integration, followed by multiplication by $z+R$, bounds the change
in squared-energy offset by
\begin{equation}\label{inf:wide-low-error}
 C C_{\mathrm{ap}}(\ell+1)/p_0
       +C C_{\mathrm{ap}}(1+C_{\mathrm{ap}})/p_0^2.
\end{equation}
For $\ell=O(\sqrt{p_0})$ this tends to zero in both stated ranges.
The offset at $p_0$ is then less than $64$, contradicting the first
conclusion of Proposition~\ref{inf:low-exclusion}. The other radial
seed states are excluded by their frequency spacing and the same
continuation estimate.
\end{proof}

\begin{lemma}[Extension of the split interval]\label{inf:split-extension}
The reduction, centre, flow, coefficient-space, material, mixed,
localization and cofactor estimates above hold on $I_{\mathrm{sp}}$
with the replacements in the following table. Every constant may
depend on $\varepsilon_{\mathrm{sp}}$; their powers of $p$ and their
order of choice are unchanged. The weighted reduction, centre and
flow statements allow real $r$ on this interval. The Euclidean
coefficient, material, mixed and closing estimates are applied at
integer block sizes $a,b$ with $a+b=2p$.
For the $\tau\le20$ counting estimate
one may use
\[
 K_{\mathrm{sp}}=\frac4{\varepsilon_{\mathrm{sp}}
                            (1-\varepsilon_{\mathrm{sp}})},\qquad
 W_{\mathrm{sp}}=40K_{\mathrm{sp}}.
\]
\begingroup\small\normalfont
\renewcommand{\arraystretch}{1.13}
\begin{longtable}{@{}>{\raggedright\arraybackslash}p{.29\textwidth}>{\raggedright\arraybackslash}p{.29\textwidth}>{\raggedright\arraybackslash}p{.35\textwidth}@{}}
\toprule Statement & Where the interval enters & Replacement\\\midrule
\endfirsthead
\toprule Statement & Where the interval enters & Replacement\\\midrule
\endhead
\ref{inf:reduction}, \ref{inf:angular-sup}, \ref{inf:flow-form} &
Positive beta parameters, endpoint orientation, axis cutoffs &
$\ab,\bb\ge\varepsilon_{\mathrm{sp}}p>1$ and $\bb\ge\ab$;
integer blocks for the Euclidean identification.\\
\ref{inf:quartic-normalization} & $pg_2>9/16$,
$\nrm{\mathsf H_2}_\infty<9/16$, $\nrm{T_\varsigma}_\infty<1$ &
$pg_2\ge2\varepsilon_{\mathrm{sp}}(1-\varepsilon_{\mathrm{sp}})$,
$pg_2<1$, $\nrm{\mathsf H_2}_\infty\le1$,
$\nrm{T_\varsigma}_\infty\le K_{\mathrm{sp}}$.\\
\ref{inf:angular-weighted} & Derivatives of $\sqrt\varsigma$,
lower bound on $\mathfrak D_j$ &
$\varsigma\ge\varepsilon_{\mathrm{sp}}^2$,
$\mathfrak D_j\ge c_{\mathrm{sp}}(p+j)^2$;
constants $C_{\mathrm{sp}}$ in the same weighted estimates.\\
\ref{inf:angular-positive} & Adjacent coefficient ratio;
the upper bound $g_2<1/p$ &
$\ab,\bb\ge16$ suffices for the same ratio bounds;
$\mathcal N_j,\mathcal E_j$ and their certificates are
independent of $\varsigma$.\\
\ref{inf:centres}, \ref{inf:centre-coefficient-defect} &
Compact parameter set and divisors $r(1-r)$ &
Use $r\in[\varepsilon_{\mathrm{sp}},1/2-\varepsilon_{\mathrm{sp}}]$;
the same finite recursion and each prescribed complex disc.\\
\ref{inf:flow-relative}--\ref{inf:flow-theorem} &
Quartic multiplier bounds and positive angular margin &
$C_{\mathrm{sp}}$ replaces numerical interval constants;
then choose $W_{\mathrm{flow}}$ and the threshold.\\
\ref{inf:quartic-count-window} & Supremum of $T_\varsigma$ &
Shell width $3K_{\mathrm{sp}}/p^2$ and window
$W_{\mathrm{sp}}$ for sufficiently large $p$.\\
\ref{inf:grid-avoidance}, \ref{inf:integer-gap} &
Grid size and spacing &
$\#\{a:2\varepsilon_{\mathrm{sp}}p\le a\le
(1-2\varepsilon_{\mathrm{sp}})p\}$ and spacing
at least $\varepsilon_{\mathrm{sp}}/p$.\\
\eqref{inf:gasper-hypothesis}, Section~\ref{inf:sec-spaces} &
Gasper margin for $\bb-\ab\le p/2$ &
Set $u=p-1$, $v=\bb-\ab\le u$; the lower margin is
$18u^3+63u^2+45u>0$.\\
\ref{inf:uniform-centre-size}, \ref{inf:relative-graph} &
The bounds for $pg_2$ and $\mathsf H_2$ &
$C_0,C_{\mathrm{rel}}$ depend on the split interval;
choose $E_{\mathrm{inv}}$, then $M$, thereafter.\\
\ref{inf:material-isomorphism}, \ref{inf:geometry-nonball} &
Normal Hessian and the physical seed coefficient &
$\nrm{D_p^{-1}}_{\Ec_1}\le C_{N,\mathrm{sp}}p^{-2}$ and
$\mathsf H_2(1)\ge c_{\mathrm{sp}}>0$;
the seed is $-\theta/(16r(1-r)p)+O_{N,\mathrm{sp}}(\tau p^{-2})$.\\
\bottomrule
\end{longtable}
\endgroup
\end{lemma}
\begin{proof}
The rational identity \eqref{inf:g2-exact} is valid for all these real
parameters. Its leading term is $pg_2\to4r(1-r)$ uniformly on the
compact interval, and its negative $\varsigma$ coefficient gives
$pg_2<1$ as before. The two endpoint values of $\mathsf H_2$ tend
uniformly to $r^2$ and $(1-r)^2$; its negative minimum is
$O(p^{-1})$. This proves the displayed replacement bounds and
$\mathsf H_2(1)\ge c_{\mathrm{sp}}$. The formulas for
$\mathfrak D_j$, $\mathfrak a_j$, and $\mathfrak b_j$ in
Lemma~\ref{inf:angular-weighted} then give every entry and derivative
estimate with $C_{\mathrm{sp}}$.

For positivity, impose $p\varepsilon_{\mathrm{sp}}\ge16$. The two
adjacent-ratio computations in Lemma~\ref{inf:angular-positive} use
only $\ab,\bb\ge16$, $j\ge7$, and $2j+p\ge78$.
Its numerators \eqref{inf:N-angular}, \eqref{inf:E-angular} and
\eqref{inf:half-diagonal} have no $\varsigma$. The last lower bound
uses $g_2<1/p$, which was retained. The same positive variation
therefore holds. In the form proof all series converge when
$C_{\mathrm{sp}}\tau/p^2$ is small; its two Schur reductions then
give the stated flow estimates with the changed constants.

For the counting row, the normalized radius displacement is at most
$10K_{\mathrm{sp}}/R^2\le3K_{\mathrm{sp}}/p^2$ for $\tau\le20$.
The min--max calculation of Lemma~\ref{inf:quartic-count-window}
then gives a containing window $40K_{\mathrm{sp}}$ for all sufficiently
large $p$ depending on $K_{\mathrm{sp}}$. For consecutive integers $a$,
integrating $|d\varsigma/da|=|r-1/2|/p$ gives spacing at least
$\varepsilon_{\mathrm{sp}}/p$. Use this spacing in the clipped-curve count.
For the one-split gap, take $\gamma=c_{\mathrm{sp}}'\tau p^{-5/3}$.
The clipped-curve count is at most
$C_{\mathrm{sp}}p^{2/3}+C_{\mathrm{sp}}c_{\mathrm{sp}}'p$.
Choose $c_{\mathrm{sp}}'>0$ so the second term is less than half
the leading grid cardinality, then increase $p$ to absorb the first.
This gives the $p^{-5/3}$ gap on the extended interval as well.

For the algebra put $u=\ab+\bb-1=p-1$ and $v=\bb-\ab$.
If $u^2-7u-24\le0$, \eqref{inf:gasper-hypothesis} holds directly.
Otherwise its left-minus-right side is minimized on $0\le v\le u$
at $v=u$, where it equals $18u^3+63u^2+45u$.
Positive products, endpoint bounds and real radial stencils follow.
The material and mixed proofs subsequently use these products and
the strong norm of the finite centre; their numerical use of $9/16$
is replaced by the first two rows. The recurrence for that centre
divides only by nonzero units, $-1/8$, $c_j$ and powers of $r(1-r)$.
These divisors stay uniformly away from zero on the compact interval.
Thus all listed estimates follow in their stated order.
\end{proof}

\begin{lemma}[Arbitrary fixed detuning bound]\label{inf:fixed-detuning-extension}
Fix $0<C_{\mathrm{ap}},T<\infty$ and the split interval above.
For half-lattice inputs with $|p-p_0|\le C_{\mathrm{ap}}/p_0$ and
$0<\tau\le T$, the centre, flow, cofactor inverse, and Newton
statements have the parameterized forms below. Fix
$5/3<q_{\mathrm{gap}}<2$; the base case $q_{\mathrm{gap}}=5/3$
is included with its original bounds. In the inverse and closing
rows assume a selected quartic gap $\kappa\tau p^{-q_{\mathrm{gap}}}$,
with fixed $\kappa>0$.
Use $\varsigma\in I_{\mathrm{sp}}$, equivalently
$r\in[\varepsilon_{\mathrm{sp}},1/2-\varepsilon_{\mathrm{sp}}]$;
material and closing statements use integer blocks.
Constants and thresholds may depend on these fixed data.
Choose $W_{\mathrm{flow}}(T)$ above the base lower bound and the
required multiples of $T$ and of the inverse angular margin.
After $C_0,C_{\mathrm{rel}}$ choose $E_{\mathrm{inv}}$, then
$W_{\mathrm{loc}}=TE_{\mathrm{inv}}$ and $M=M_{TE_{\mathrm{inv}}}$,
then separation and the head cutoff, and finally the centre order.
The last column identifies the proof line permitting each substitution.

\begingroup\small\normalfont
\def\UrlBreaks{\do\:\do\-}
\setlength{\tabcolsep}{4pt}
\renewcommand{\arraystretch}{1.16}
\begin{longtable}{@{}>{\raggedright\arraybackslash}p{.23\textwidth}>{\raggedright\arraybackslash}p{.16\textwidth}>{\raggedright\arraybackslash}p{.29\textwidth}>{\raggedright\arraybackslash}p{.25\textwidth}@{}}
\toprule Base label & Parameter replaced & Parameterized bound & Proof line permitting replacement\\\midrule
\endfirsthead
\toprule Base label & Parameter replaced & Parameterized bound & Proof line permitting replacement\\\midrule
\endhead
\ref{inf:detuning}, \nolinkurl{inf:detuning};
\ref{inf:low-exclusion}, \nolinkurl{inf:low-exclusion}
& Scaled distance $31/10$, base detuning bound
& Use $|p-p_0|\le C_{\mathrm{ap}}/p_0$ and
$0<\tau\le T$; the radius and detuning comparison are
\eqref{inf:log-detuning-relation}. Low exclusion retains its
original energy bounds.
& The exact quotient \eqref{inf:trace-ratio-exact} and the real
zero derivative comparison hold on this fixed window;
\ref{inf:wide-low-transfer} gives the continuation error
\eqref{inf:wide-low-error}.\\
\ref{inf:centres}, \nolinkurl{inf:centres};
\ref{inf:centre-coefficient-defect}, \nolinkurl{inf:centre-coefficient-defect}
& $|\theta|\le20$, base split interval
& \eqref{inf:centre-complex-defect} and
\eqref{inf:centre-coeff-defect-bound} are
$C_{N,T,\mathrm{sp}}\theta^2p^{-N}$;
\eqref{inf:centre-profile} has physical error
$C_{N,T,\mathrm{sp}}\tau p^{-3}$.
& In \eqref{inf:rho-equation} the $Z$ derivative at $(\eps,Z)=(0,4)$
is $8$ for every $\theta$; the recursion divides by
units, $-1/8$, $c_j$ and $r(1-r)$.\\
\ref{inf:flow-relative}, \nolinkurl{inf:flow-relative};
\ref{inf:flow-complement}, \nolinkurl{inf:flow-complement}
& $\tau\le20$, absolute flow window
& Relative bounds $C_{\mathrm{sp}}\tau p^{-2}$ and
$C_{\mathrm{sp}}\tau^2p^{-3}$;
\eqref{inf:flow-schur-remainder} is
$C_{\mathrm{sp}}\tau^2(W_{\mathrm{flow}}(T)^{-1}+p^{-1})$.
& The series is $\sum h^6(C_{\mathrm{sp}}\tau/p^2)^h$;
the complement ratio is $C_{\mathrm{sp}}\tau/W_{\mathrm{flow}}(T)$.
The derivative adds $C_{\mathrm{sp}}\tau^3/W_{\mathrm{flow}}(T)^2$.\\
\ref{inf:flow-low-block}, \nolinkurl{inf:flow-low-block};
\ref{inf:ordered-eigenvalues}, \nolinkurl{inf:ordered-eigenvalues};
\ref{inf:flow-theorem}, \nolinkurl{inf:flow-theorem}
& $\tau\le20$, $\varsigma_0\in(1/64,1/16)$
& Use $\varsigma_0\in\operatorname{int}I_{\mathrm{sp}}$;
the low inverse is $C_{T,\mathrm{sp}}/(p\tau)$.
For $|\lambda|<\tau/(16p)$ the signed slope is
at least $c_{T,\mathrm{sp}}\tau/p$.
& The proof of \ref{inf:flow-low-block} adds the first-order term
to give seed diagonal $\theta/(2p)+O_T(\theta^2p^{-3})$;
\eqref{inf:low-offsets} keeps the index-five value
$-72+O_T(p^{-1})$. Final norm recovery in
\ref{inf:flow-theorem} costs $C_{T,\mathrm{sp}}\sqrt p$.\\
\ref{inf:quartic-count-window}, \nolinkurl{inf:quartic-count-window}
& $|\theta|\le20$, $W_{\mathrm{count}}=40$
& Shell width $C_{\mathrm{sp}}\tau/p^2$ and containing
count window $C_{\mathrm{sp}}(1+T)$.
& \eqref{inf:quartic-domain} gives displacement
$\tau\nrm{T_\varsigma}_\infty/(2R^2)$;
indexwise shell bounds multiply $O(p^2)$ eigenvalues by
$1+O_{\mathrm{sp}}(\tau/p^2)$.\\
\ref{inf:grid-avoidance}, \nolinkurl{inf:grid-avoidance};
\ref{inf:integer-gap}, \nolinkurl{inf:integer-gap}
& Base split grid and $\tau\le20$
& The interior grid has $\asymp_{\mathrm{sp}}p$ points and spacing
at least $\varepsilon_{\mathrm{sp}}/p$; one split has gap
$c_{T,\mathrm{sp}}\tau p^{-5/3}$.
& The clipped-curve count is
$C_{T,\mathrm{sp}}p^{2/3}(1+C_{T,\mathrm{sp}}p\gamma/\mu)$,
with $\mu=c_{T,\mathrm{sp}}\tau/p$;
choose the gap constant before the threshold.\\
\ref{inf:uniform-centre-size}, \nolinkurl{inf:uniform-centre-size};
\ref{inf:relative-graph}, \nolinkurl{inf:relative-graph}
& Centre-size and relative constants on the base compact set
& $\delta\le C_{0,\mathrm{sp}}\tau p^{-2}$;
\eqref{inf:relative-scalar} and \eqref{inf:relative-two-norms}
are $C_{\mathrm{rel},\mathrm{sp}}\tau p^{-2}$.
These constants precede $N$.
& The leading quartic size is $C_{\mathrm{sp}}\tau p^{-2}$;
the finite-order remainder $C_{N,T,\mathrm{sp}}\tau p^{-4}$
is absorbed by the threshold. The scalar graph line is
$C(\delta+p\delta^2)$.\\
\ref{inf:centre-source-norm}, \nolinkurl{inf:centre-source-norm};
\ref{inf:residual-estimates}, \nolinkurl{inf:residual-estimates}
& Fixed-centre constants for $\tau\le20$
& \eqref{inf:centre-norm-bounds} retains powers $p^{k+1}$ and
$p^{k+5}$ with $C_{N,k,T,\mathrm{sp}}$.
Residual and material defect bounds are
$C_{N,T,\mathrm{sp}}\theta^2p^{2-N}$ and
$C_{N,T,\mathrm{sp}}\theta^2p^{3-N}$.
& The known-mode head uses the finite factor $e^{C_N\sqrt T}$;
the factorial radial tail stays geometric. The two-layer lift
and \eqref{inf:Delta-moment} give the residual powers.\\
\ref{inf:material-isomorphism}, \nolinkurl{inf:material-isomorphism};
\eqref{inf:p6-second}, \nolinkurl{inf:p6-second}
& Fixed-centre constants for $\tau\le20$
& $\nrm{\Qmat^{-1}}\le C_{N,T,\mathrm{sp}}p^6$ and
$\sup\nrm{D^2\Fnl}\le C_{N,T,\mathrm{sp}}p^6$
on \eqref{inf:nonlinear-working-region}.
& \eqref{inf:material-reconstruction} uses
$\nrm{D_p^{-1}}\le C_{N,T,\mathrm{sp}}p^{-2}$ and
$\nrm{\Delta T_\eta}\le C_{N,T,\mathrm{sp}}p^8\nrm\eta$;
\eqref{inf:nonlinear-second} costs $Cp(1+\nrm g)$.\\
\eqref{inf:inverse-window}, \nolinkurl{inf:inverse-window};
\ref{inf:inverse-complement}, \nolinkurl{inf:inverse-complement};
\ref{inf:inverse-Schur}, \nolinkurl{inf:inverse-Schur}
& Containing window $20E_{\mathrm{inv}}$
& The projection stays
$\one_{[-E_{\mathrm{inv}}\tau,E_{\mathrm{inv}}\tau]}(H_0-R^2)$;
its fixed superset has width $TE_{\mathrm{inv}}$.
The Schur norm is $C_{E_{\mathrm{inv}},T,\mathrm{sp}}\tau$.
& \eqref{inf:ball-complement-bound} and
\eqref{inf:relative-two-norms} give complementary absorption;
the four feedback factors give $C\tau/E_{\mathrm{inv}}$.\\
\ref{inf:cluster-inverse}, \nolinkurl{inf:cluster-inverse}
& $\tau\le20$, gap exponent
& For every $\tau>0$, gap $c\tau p^{-q_{\mathrm{gap}}}$ and
$\sigma D>q_{\mathrm{gap}}M+q_{\mathrm{gap}}+2$ give
$\nrm{F^{-1}}_1\le C\tau^{-1}p^{q_{\mathrm{gap}}M}$.
& The cofactor line is $m!(C\tau)^{m-1}(2/\gamma_F)^m$;
$\tau$ cancels in both error absorptions.\\
\ref{inf:effective-inverse}, \nolinkurl{inf:effective-inverse}
& $M_{20E_{\mathrm{inv}}}$, $5M/3$, $D_{\mathrm{sep}}>5M/3+5$
& $M=M_{TE_{\mathrm{inv}}}$,
$\nrm{\Feff^{-1}}_1\le C\tau^{-1}p^{q_{\mathrm{gap}}M}$;
choose $\sigma D_{\mathrm{sep}}>q_{\mathrm{gap}}M+q_{\mathrm{gap}}+2$.
& \ref{inf:clusters} applies to the fixed superset;
deleting inter-cluster entries costs
$C\tau p^{-\sigma D_{\mathrm{sep}}}$, followed by the cofactor row.\\
\ref{inf:mixed-inverse}, \nolinkurl{inf:mixed-inverse}
& Exponent $2+5M/3$
& \eqref{inf:mixed-inverse-bound} has right side
$C\tau^{-1}p^{2+q_{\mathrm{gap}}M}$.
& The block solution costs $C/E_{\mathrm{inv}}$ off diagonal,
$C\tau$ into the window and $Cp^2$ for $K_0\Pi$;
density graph multipliers are bounded.\\
\ref{inf:coefficient-inverse}, \nolinkurl{inf:coefficient-inverse}
& $\tau\le20$, exponent $5/2+5M/3$
& $\nrm{J_D^{-1}}_\Cc\le C\tau^{-1}p^{5/2+q_{\mathrm{gap}}M}
\le Cp^{2M+3}/\gamma_p$, with
$\gamma_p=c_{\mathrm{gap}}\tau p^{-q_{\mathrm{gap}}}$.
& \eqref{inf:head-conversion} costs $\sqrt p$;
\eqref{inf:relative-scalar} gives \eqref{inf:head-tail-hypotheses}.
Compare $5/2+q_{\mathrm{gap}}(M-1)\le2M+3$.\\
\ref{inf:centre-gap-transfer}, \nolinkurl{inf:centre-gap-transfer}
& Gap exponent $5/3$, $\tau\le20$
& Quartic gap $c_{\mathrm q}\tau p^{-q_{\mathrm{gap}}}$
gives centre gap $(c_{\mathrm q}/2)\tau p^{-q_{\mathrm{gap}}}$.
& The min--max displacement is $C_{N,T,\mathrm{sp}}\tau p^{-2}$;
require $p^{2-q_{\mathrm{gap}}}>2C_{N,T,\mathrm{sp}}/c_{\mathrm q}$.\\
\ref{inf:newton-threshold}, \nolinkurl{inf:newton-threshold};
\ref{inf:Newton-zero}, \nolinkurl{inf:Newton-zero}
& $\tau\le20$, gap exponent $5/3$
& Inverse hypothesis
$C_Np^{b_{\mathrm{inv}}}/(c_{\mathrm{gap}}\tau p^{-q_{\mathrm{gap}}})$;
$B=b_{\mathrm{inv}}+m_{\mathrm{mat}}+q_{\mathrm{gap}}$,
$L=B+d_{\mathrm{nl}}+4$, $N\ge\lceil2B+d_{\mathrm{nl}}+8\rceil$.
& \eqref{inf:newton-gate-one}--\eqref{inf:newton-gate-three} give
$C_Np^{B+L+2-N}$, $C_N\tau p^{B+3-N}$ and
$C_Np^{B+d_{\mathrm{nl}}-L}$; only the second uses $\tau\le T$.\\
\ref{inf:geometry-convexity}, \nolinkurl{inf:geometry-convexity};
\ref{inf:geometry-nonball}, \nolinkurl{inf:geometry-nonball}
& Fixed-centre comparison constants
& Normalized $C^2$ size
$C_{N,T,\mathrm{sp}}\tau p^{-2}+O(\tau p^{-L})$;
physical seed $-\theta/(16r(1-r)p)$ with error
$C_{N,T,\mathrm{sp}}\tau p^{-2}+O(\tau p^{1-L})$.
& \eqref{inf:curvature-numerator} applies when the shape tends
to zero in $C^2$; dividing the witness error by $\tau/p$
makes it tend to zero.\\
\bottomrule
\end{longtable}
\endgroup
In the Newton row take $b_{\mathrm{inv}}=2M+3$ and
$m_{\mathrm{mat}}=d_{\mathrm{nl}}=6$. Thus
$B=B_{\mathrm{bd}}=2M+9+q_{\mathrm{gap}}$,
$L=L_{\mathrm{bd}}=B_{\mathrm{bd}}+10$ and $N=4M+40$;
the three powers are those in \eqref{inf:bounded-gates}.
Section~\ref{inf:sec-constants} records the special case
$T=20$, $q_{\mathrm{gap}}=5/3$.
These fixed-detuning families use the cofactor estimate
with fixed local cluster cardinality.
\end{lemma}
\begin{proof}
Lemma~\ref{inf:wide-low-transfer} supplies the low-state exclusion.
The real root comparison and the exact ratio
\eqref{inf:trace-ratio-exact} give $R=2p+3+O_{C_{\mathrm{ap}}}(p^{-1})$.
For the centre choose a fixed complex $\theta$ disc containing
$[-T,T]$ before its $\eps$ disc. Its radius equation has derivative
eight at $(\eps,Z)=(0,4)$ for every $\theta$ on this fixed disc.
The finite recursion has the same nonzero leading divisors, and its
two defects have the same analytic factor $\theta^2$. Compactness
therefore proves all the stated fixed-order centre bounds with
$C_{N,T,\mathrm{sp}}$.

In Lemma~\ref{inf:flow-relative} the dominating series is
$\sum h^6(C_{\mathrm{sp}}\tau/p^2)^h$. Choose a fixed
$W_{\mathrm{flow}}$ larger than the required multiples of $T$ and
the inverse angular margin. The complementary series ratio is
$C_{\mathrm{sp}}\tau/W_{\mathrm{flow}}$, and its derivative remainder
is bounded by $C_{\mathrm{sp}}\tau^2(W_{\mathrm{flow}}^{-1}+p^{-1})$.
The extra term $C\tau^3/W_{\mathrm{flow}}^2$ is absorbed by this
bound. The seed diagonal is still
$\theta/(2p)+O_T(\theta^2p^{-3})$, the index-five diagonal is
$-72+O_T(p^{-1})$, and the weighted low inverse is $C_T/(p\tau)$.
Its two exterior remainders retain their factors $\tau^2$.
Choose the window and then $p$ so the errors are below the positive
angular margin. The last step of Theorem~\ref{inf:flow-theorem}
therefore gives $c_{T,\mathrm{sp}}\tau/p$ in its same narrow band.
The shell count uses a fixed window $C_{\mathrm{sp}}(1+T)$.

The leading quartic strong norm is bounded by
$C_{\mathrm{sp}}\tau p^{-2}$ independently of the centre order.
For each fixed $N$ its remainder is
$C_{N,T,\mathrm{sp}}\tau p^{-4}$. Increasing only the threshold
absorbs this remainder into the leading bound, as in
Lemma~\ref{inf:uniform-centre-size}. Thus $C_0,C_{\mathrm{rel}}$
are fixed before $N$. For the coefficient inverse choose
$E_{\mathrm{inv}}$ from this relative graph constant. The
containing window is now $TE_{\mathrm{inv}}$, so
$M=M_{TE_{\mathrm{inv}}}$ is fixed. The determinant proof of
Lemma~\ref{inf:cluster-inverse} retains exactly one power $\tau^{-1}$.
Choose separation and head cutoff after $M$, and the finite centre
order after those choices. The material and residual proofs use
only the fixed-centre norms and the unchanged clamped estimates.
Their constants may depend on $T$, while their exponents remain six.
The three Newton gate estimates and the geometric comparisons
retain their powers of $\tau$; replacing its upper bound twenty by
$T$ changes only the threshold.
\end{proof}

\begin{lemma}[All but a sublinear number of interior splits]
\label{inf:many-splits}
Fix $C_{\mathrm{ap}},T$, $0<\eta_{\mathrm{cnt}}<1/3$, and
$q_{\mathrm{gap}}=2-\eta_{\mathrm{cnt}}$. At every sufficiently large input of
Lemma~\ref{inf:fixed-detuning-extension}, all but
$O_{C_{\mathrm{ap}},T,\mathrm{sp},\eta_{\mathrm{cnt}}}
 (p^{2/3+\eta_{\mathrm{cnt}}})$ integers
\[
 2\varepsilon_{\mathrm{sp}}p\le a\le(1-2\varepsilon_{\mathrm{sp}})p
\]
give an actual quartic invariant gap at least
$\kappa\tau p^{-q_{\mathrm{gap}}}$, where $\kappa>0$ depends only
on the fixed data.
\end{lemma}
\begin{proof}
The enlarged shell window gives at most $C_{T,\mathrm{sp}}p^{2/3}$
possible ordered curves. Their clipped signed functions are
nondecreasing by the absolute-continuity argument of
Lemma~\ref{inf:grid-avoidance}, with slope at least
$\mu=c_{T,\mathrm{sp}}\tau/p$ inside the narrow flow band.
Since the grid spacing is at least $\varepsilon_{\mathrm{sp}}/p$,
one curve excludes at most
\[
 1+\frac{2p\kappa\tau p^{-q_{\mathrm{gap}}}}
          {\varepsilon_{\mathrm{sp}}\mu}
 \le 1+C p^{2-q_{\mathrm{gap}}}
\]
points. The gap lies inside the flow band eventually because
$q_{\mathrm{gap}}>1$. Summing the exclusions gives
$C(p^{2/3}+p^{8/3-q_{\mathrm{gap}}})\le C'p^{2/3+\eta_{\mathrm{cnt}}}$.
One can exclude twice the stated radius first, which also covers
equality at an endpoint. This proof uses $2p\in\Z$ only to obtain
integer second block sizes.
\end{proof}

\begin{proposition}[Cofactor closure]
\label{inf:bounded-assembly}
Under the hypotheses of Lemma~\ref{inf:many-splits}, every retained
split carries an analytic strictly convex non-ball domain and a
solution of \eqref{inf:main-equations} with its prescribed block symmetry.
The threshold in $p$ follows these choices.
First fix $C_{\mathrm{ap}},T,\varepsilon_{\mathrm{sp}},q_{\mathrm{gap}}$.
Next fix $C_0,C_{\mathrm{rel}},E_{\mathrm{inv}}$ and
$M=M_{TE_{\mathrm{inv}}}$, then $D_{\mathrm{sep}}$ and
$K_{\mathrm{head}}$. Finally set
\[
 b_{\mathrm{inv}}=2M+3,\qquad N=4M+40,\qquad
 B_{\mathrm{bd}}=2M+9+q_{\mathrm{gap}},\qquad
 L_{\mathrm{bd}}=B_{\mathrm{bd}}+10.
\]
The closed Newton ball has radius $\tau p^{-L_{\mathrm{bd}}}$.
\end{proposition}
\begin{proof}
Fix the data in the order stated, choosing separation with
$\sigma D_{\mathrm{sep}}>q_{\mathrm{gap}}M+q_{\mathrm{gap}}+2$
and the head cutoff required by Lemma~\ref{inf:head-tail}.
Use the parameterized statements of
Lemma~\ref{inf:fixed-detuning-extension} throughout this chain.
The centre-to-quartic normalized displacement is
$C_{N,T,\mathrm{sp}}\tau p^{-4}$. The indexwise min--max proof of
Proposition~\ref{inf:centre-gap-transfer} changes energies by
$C_{N,T,\mathrm{sp}}\tau p^{-2}=o(\tau p^{-q_{\mathrm{gap}}})$.
Thus the centre has half the selected gap.

Use Lemma~\ref{inf:cluster-inverse} with this gap exponent.
The Schur reconstruction and radial conversion in
Section~\ref{inf:sec-inverse} give
\[
 \nrm{J_D^{-1}}\le C\tau^{-1}p^{5/2+q_{\mathrm{gap}}M}
 \le C p^{2M+3}/(\tau p^{-q_{\mathrm{gap}}}),
\]
since $5/2+q_{\mathrm{gap}}(M-1)\le2M+3$.
The separation inequality makes the off-cluster restoration small.
The scalar estimate \eqref{inf:relative-scalar} supplies
\eqref{inf:head-tail-hypotheses} for fixed $T$ and sufficiently
large $p$, as in the proof of Theorem~\ref{inf:coefficient-inverse}.
The material inverse and nonlinear bound still have exponent six,
so $\AN=\Qmat^{-1}J_D^{-1}$ has norm at most
$C_N\tau^{-1}p^{B_{\mathrm{bd}}}$.

Apply the contraction proof of Theorem~\ref{inf:newton-threshold}
with these explicit exponents. The three gate powers are
\begin{equation}\label{inf:bounded-gates}
 \begin{split}
 B_{\mathrm{bd}}+L_{\mathrm{bd}}+2-N&=2q_{\mathrm{gap}}-10,\\
 B_{\mathrm{bd}}+3-N&=-2M-28+q_{\mathrm{gap}},\\
 B_{\mathrm{bd}}+6-L_{\mathrm{bd}}&=-4.
 \end{split}
\end{equation}
The second carries the additional factor $\tau\le T$; all three
tend to zero uniformly in $0<\tau\le T$. The working ball lies in
\eqref{inf:nonlinear-working-region} for large $p$. The actual inverse
identities in the same spaces give the exact real zero. The geometric
and classical-solution arguments of Section~\ref{inf:sec-geometry}
use the replacement constants in Lemma~\ref{inf:split-extension}.
Its normalized $C^2$ size is
$O_{N,T,\mathrm{sp}}(\tau p^{-2})+O(\tau p^{-L_{\mathrm{bd}}})\to0$,
so Lemma~\ref{inf:geometry-convexity} gives strict convexity.
The leading physical quartic coefficient has size comparable to
$\tau/p$; both the centre remainder and the physical Newton
correction are smaller after division by that quantity.
This proves non-sphericity as well as strict convexity, and completes
the cofactor chain uniformly for $0<\tau\le T$.
\end{proof}

\begin{lemma}[Orthogonal symmetry and similarity classes]
\label{inf:similarity-classes}
For a differentiable nonconstant radial graph about zero invariant
under $O(a)\times O(n-a)$, its orthogonal symmetry Lie algebra is
$\mathfrak{so}(a)\oplus\mathfrak{so}(n-a)$.
Consequently constructed domains at distinct splits below $n/2$
are pairwise non-similar. For a dimension family, define
$\mathcal N_{\mathrm{sim}}(n)$ as the number of integers $0<a<n/2$
for which the construction gives a domain on some fixed compact
interior split interval, with its associated fixed centre order.
Choosing any one such domain at each split gives this many
similarity classes, independently of the choices.
Suppose that, for every fixed $0<\varepsilon_{\mathrm{sp}}<1/4$
and $0<\eta_{\mathrm{cnt}}<1/3$, all but
$O(n^{2/3+\eta_{\mathrm{cnt}}})$ integer splits in
$[\varepsilon_{\mathrm{sp}}n,(1/2-\varepsilon_{\mathrm{sp}})n]$
give domains. Then along that family's dimensions,
\[
 \mathcal N_{\mathrm{sim}}(n)\ge
 (1/2-2\varepsilon_{\mathrm{sp}})n-Cn^{2/3+\eta_{\mathrm{cnt}}},\qquad
 \liminf_{n\to\infty}\frac{\mathcal N_{\mathrm{sim}}(n)}n\ge\frac12.
\]
The limit inferior is taken through that family's dimensions.
A countable exhaustion by fixed compact split intervals suffices
in the definition of $\mathcal N_{\mathrm{sim}}$.
\end{lemma}
\begin{proof}
Write the radius on the sphere as $F(x)$, $x=|X|^2$.
There is $x_0\in(0,1)$ with $F'(x_0)\ne0$.
A skew matrix has diagonal blocks in the two block Lie algebras and
off-diagonal blocks $M,-M^t$. Its infinitesimal action on $F$ is
$2F'(x)X^tMY$. Vanishing for all $X,Y$ with squared norms
$x_0,1-x_0$ forces $M=0$, by varying their directions. The block
diagonal matrices are symmetries, proving the equality.

All constructed domains have barycentre zero. A similarity between
two has zero translation and conjugates their orthogonal symmetry
Lie algebras. The dimension of the latter is
$\{a(a-1)+(n-a)(n-a-1)\}/2$, whose change from $a$ to $a+1$ is
$2a+1-n<0$ on the selected range. Thus different splits give different
similarity classes. The assumed exceptional-split bound and integer
counting give the displayed lower bound. Fixing any $\varepsilon_{\mathrm{sp}}>0$
and then taking the lower limit proves a bound $1/2-2\varepsilon_{\mathrm{sp}}$.
Taking the supremum over these fixed choices proves the last assertion;
their thresholds need not be uniform as $\varepsilon_{\mathrm{sp}}\downarrow0$.
\end{proof}

\begin{theorem}[Dimensions in prescribed residue classes]
\label{inf:residue-arithmetic}
For every positive integer $L_{\mathrm{ap}}$ and every integer
$b_{\mathrm{ap}}$, there is an unbounded selection of dimensions
$n\equiv b_{\mathrm{ap}}\pmod{4L_{\mathrm{ap}}}$ with
\[
 |n/2-\widehat p_m|<32L_{\mathrm{ap}}^2/\widehat p_m,
 \qquad 0<\tau<256L_{\mathrm{ap}}^2
\]
eventually. Every sufficiently large selected dimension has the
domains and multiplicity just proved. The threshold may depend on
the fixed modulus, residue, and split/counting parameters.
In particular every residue modulo $L_{\mathrm{ap}}$ occurs
infinitely often. For even $L_{\mathrm{ap}}$ the residue determines
the parity; for odd $L_{\mathrm{ap}}$ both parities occur infinitely
often in each residue class modulo $L_{\mathrm{ap}}$.
\end{theorem}
\begin{proof}
Set $M_{\mathrm{ap}}=4L_{\mathrm{ap}}$. Consecutive convergent vectors
of $\alpha_{\mathrm{ar}}$ form an invertible matrix modulo
$M_{\mathrm{ap}}$. Choose representatives
$A,B\in\{0,\ldots,M_{\mathrm{ap}}-1\}$ satisfying
\[
 P=A P_k^{\mathrm{cf}}+B P_{k-1}^{\mathrm{cf}}
       \equiv b_{\mathrm{ap}}+3\pmod{M_{\mathrm{ap}}},\qquad
 Q=A Q_k^{\mathrm{cf}}+B Q_{k-1}^{\mathrm{cf}}
       \equiv1\pmod{M_{\mathrm{ap}}}.
\]
They cannot both vanish, so $Q\ge Q_{k-1}^{\mathrm{cf}}>0$.
In \eqref{inf:cf-error}, for $T_k^{\mathrm{cf}}\ge A/B$ the expression is at
most $B(A+B)$; when $1<T_k^{\mathrm{cf}}<A/B$ its maximum is attained at
$T_k^{\mathrm{cf}}=1$ and then is at most
$|A-B|\max\{A,(A+B)/2\}$. The cases $A=0$ or $B=0$ follow
directly. Both bounds are at most $2(M_{\mathrm{ap}}-1)^2$.
Set $m=(Q-1)/4$, $n=P-3$, and $p=n/2$. The crossing expansion gives
\[
 \limsup\widehat p_m|p-\widehat p_m|
 \le\frac{\alpha_{\mathrm{ar}}}{2}(M_{\mathrm{ap}}-1)^2
       +d_{\mathrm{ar}}<32L_{\mathrm{ap}}^2.
\]
For the last strict inequality use
$\alpha_{\mathrm{ar}}<363/157$ and $d_{\mathrm{ar}}<1$.
The dimensions tend to infinity and have the prescribed residue.
Equation~\eqref{inf:log-detuning-relation}, with
$C_{\mathrm{ap}}=32L_{\mathrm{ap}}^2$, gives eventually
$\tau<(16\delta_0+o(1))32L_{\mathrm{ap}}^2<256L_{\mathrm{ap}}^2$;
non-coincidence gives $\tau>0$. Apply
Lemma~\ref{inf:fixed-detuning-extension} and
Proposition~\ref{inf:bounded-assembly} with these fixed constants.
Taking $\varepsilon_{\mathrm{sp}}=1/8$, the sublinear exclusions
leave a split in $r\in[1/4,3/8]$ as well.
For a residue modulo an odd $L_{\mathrm{ap}}$, its lifts
$b_{\mathrm{ap}}$ and $b_{\mathrm{ap}}+L_{\mathrm{ap}}$ have opposite
parities, giving the final assertion.
\end{proof}

\begin{corollary}[Almost every fixed split ratio]\label{inf:fixed-ratio}
Use the explicit continued-fraction selections in
Definition~\ref{inf:integer-sequence}, Lemma~\ref{inf:shifted-arithmetic}, and
Theorem~\ref{inf:residue-arithmetic}, retaining their convergent
index $k$. There is a set $\mathcal R_{\mathrm{split}}\subset(0,1/2)$
of full Lebesgue measure such that, for every
$r_0\in\mathcal R_{\mathrm{split}}$, each of these countably many
selected dimension families $(n_k)$ eventually has a constructed
Schiffer domain with split
\[
 a_k=\lfloor n_k r_0\rfloor.
\]
The starting index may depend on $r_0$ and the selected family.
\end{corollary}
\begin{proof}
Fix one selected family and a compact ratio interval inside
$(0,1/2)$. Choose a slightly larger compact interval for
Lemma~\ref{inf:many-splits}; rounding changes the ratio by at most
$1/n_k$, so it remains in that interval eventually.
Take $\eta_{\mathrm{cnt}}=1/6$. At dimension $n_k$ at most
$C n_k^{5/6}$ splits are excluded, with $C$ depending on this
family's fixed detuning bound and the compact interval.
Each integer split corresponds under $a=\lfloor n_k r\rfloor$
to an interval of $r$ of length $1/n_k$. The set of ratios giving
an excluded split therefore has measure at most $C n_k^{-1/6}$.

The selected denominator is at least $Q^{\mathrm{cf}}_{k-3}$
for the first two selections and at least $Q^{\mathrm{cf}}_{k-1}$
for the residue-class selection. Continued-fraction denominators
satisfy $Q^{\mathrm{cf}}_{k+2}\ge Q^{\mathrm{cf}}_{k+1}+Q^{\mathrm{cf}}_k
\ge2Q^{\mathrm{cf}}_k$. Since the selected dimensions are asymptotic
to $\alpha_{\mathrm{ar}}$ times their selected denominators,
$n_k\ge c\,2^{k/2}$ eventually. Hence $\sum_k n_k^{-1/6}<\infty$.
The measure of the union of the excluded-ratio events after index
$K$ is bounded by the tail of this convergent series and tends
to zero. Almost every ratio has only finitely many exclusions.

Take a countable compact exhaustion of $(0,1/2)$ and then the
intersection of these full-measure sets over the two primary
selections and all fixed moduli and residues. Their exceptional
sets have total measure zero. Proposition~\ref{inf:bounded-assembly}
gives the domain at every eventual retained split, proving the
assertion with the stated order of quantifiers.
\end{proof}
